% !TEX root = ../../main.tex

\renewcommand{\thesection}{1.5}

\section{因式分解定理}\label{sec:1.5}

\begin{ex}{2}\label{ex:1.5.2}
    设 \(1\) 是 \(ax^4+bx^2+1\) 的重根, 求 \(a,b\).
\end{ex}

\begin{solution}
    由 \(1\) 是 \(ax^4+bx^2+1\) 的重根, 有
    \[
    f\p{1}=0,\quad f'\p{1}=0.
    \]
    即
    \begin{gather*}
        a+b+1=0, \\
        4a+2b=0.
    \end{gather*}
    解得
    \[
    a=1,\quad b=-2.\qedhere
    \]
\end{solution}

\begin{ex}{4}\label{ex:1.5.4}
    证明多项式 \(f\) 是某个既约多项式的幂当且仅当有既约多项式 \(p\) 使得 \(f\mid p^m\),
    其中 \(m\) 为正整数.
\end{ex}

\begin{proof}
    充分性. 下证 \(f\) 为 \(p\) 的幂.

    假设 \(f\) 有不与 \(p\) 相伴的因式 \(q\) 满足 \(\deg q\geq1\), 有
    \begin{equation}
        \p{p,q}=1.
        \label{eq:1.5.4-Coprimability}
    \end{equation}
    由
    \[
    f\mid p^m,
    \]
    有
    \begin{equation}
        q\mid p^m.
        \label{eq:1.5.4-Divisibility}
    \end{equation}
    式~\eqref{eq:1.5.4-Coprimability} 与式~\eqref{eq:1.5.4-Divisibility} 矛盾,
    故 \(f\) 的所有因式都与 \(p\) 相伴, 即 \(f\) 为 \(p\) 的幂.

    必要性. 设
    \[
    f=at^k,
    \]
    其中 \(a\) 为非零常数, \(t\) 为既约多项式, \(k\) 为非负整数.

    令 \(p=t,m=k\), 则
    \[
    f=ap^m.
    \]
    由于 \(\p{f,a}=1\), 有
    \[
    f\mid p^m.\qedhere
    \]
\end{proof}